\scanpage{282}
\mspage{23}

\uline{une seule \unsure{flèche|flèches}} allant d'un objet donné vers un autre. En d'autres termes, dans de telles catégories tous les diagrammes sont commutatifs.

Bien sûr, on voit aussi que si on a une loi de formation \inserted{$L$} de chemins de type $\tau'$ à partir de chemins de type $\tau$ (avec la condition de fonctorialité p.r. à la cat. $X$ où on prend les chemins), cette loi est nécessairement \uline{fonctorielle par rapport aux chemins} : si on \deleted{\unsure{se do|se de}} a un hom
\diagram{\begin{tikzcd}[ampersand replacement=\&]c \arrow[r,"{\unsure{u|v}}"] \&c' \&\text{dans}\quad\underline{\mathrm{Ch}}_\tau(X)\end{tikzcd}}
alors la loi \inserted{$L$} ``se prolonge'' à de \unsure{tels|telles} $u$ pour donner
\diagram{\begin{tikzcd}[ampersand replacement=\&]L(c) \arrow[r,"{L(u)}"] \&L(c')\end{tikzcd}},
avec les relations habituelles
\[ L(\mathrm{id}_c)=\mathrm{id}_{L(c)},\quad L(u'u)=L(u')L(u). \]

\unsure{Cela|Celà} résulte de la description de $L$ comme une opération $L_\varphi$ en termes d'un foncteur bien déterminé
\diagram{\begin{tikzcd}[ampersand replacement=\&] I_{\tau'}\arrow[r,"{\varphi}"] \& I_\tau\end{tikzcd}},
lequel à son tour donne pour toute $X$ un \uline{foncteur}
\[ \left\{\begin{aligned}
\underline{\mathrm{Ch}}_\tau(X)&\longrightarrow\underline{\mathrm{Ch}}_{\tau'}(X)\\
c&\longmapsto c\circ\varphi
\end{aligned}\right. \]
plus précis que simplement une \uline{application} $\mathrm{Ch}_\tau(X)\longrightarrow\mathrm{Ch}_{\tau'}(X)$ entre les ens. d'objets correspondants\unsure{--|.}
% archive sheet number: 281
